% !TEX program = xelatex
\documentclass[11pt, UTF8, a4paper]{ctexart} % 如果内容很多，可以改为 ctexrep

% ==================== 宏包引入 ====================
\usepackage[margin=2.5cm]{geometry} % 页面边距
\usepackage{graphicx} % 插入图片
\usepackage{amsmath, amssymb} % 数学公式（运动学、PID计算必备）
\usepackage{hyperref} % 超链接与书签
\usepackage[table]{xcolor} % 颜色支持
\usepackage{listings} % 代码块排版
\usepackage{tabularx}
\usepackage[most]{tcolorbox} % 提示框（用于写“踩坑提示”）
\usepackage{array}
\usepackage{booktabs}
\usepackage{twemojis}

% ==================== 样式设置 ====================
% 超链接样式
\hypersetup{
    colorlinks=true,
    linkcolor=blue,
    filecolor=magenta,      
    urlcolor=cyan,
    pdftitle={RoboCup 3D 仿真组入门指南},
}

% 代码块样式设置 (适合 C++ 和 Bash 命令)
\lstset{
    backgroundcolor=\color{gray!10},   
    commentstyle=\color{green!50!black},
    keywordstyle=\color{blue},
    stringstyle=\color{orange},
    basicstyle=\ttfamily\small,
    breakatwhitespace=false,         
    breaklines=true,                 
    captionpos=b,                    
    keepspaces=true,                 
    numbers=left,                    
    numbersep=5pt,                  
    showspaces=false,                
    showstringspaces=false,
    showtabs=false,                  
    tabsize=4
}

% 自定义提示框 (用于避坑指南)
\newtcolorbox{tipbox}[1][]{
    colback=blue!5!white,
    colframe=blue!75!black,
    fonttitle=\bfseries,
    title=💡 提示/避坑: #1
}

% ==================== 文档信息 ====================
\title{\textbf{RoboCup 3D 仿真组入门指南}}
\author{你的名字/团队名}
\date{\today}

% ==================== 正文开始 ====================
\begin{document}

\maketitle
\tableofcontents % 生成目录
\newpage

\section{实现linux环境}
实现linux环境有三种方式虚拟机、WSL2和双系统。

% ================= linux环境区别表格 =================
\begin{table}[htbp]
    \centering
    
    % 定义颜色 (Color Definitions)
    \definecolor{tableheader}{RGB}{52, 73, 94}
    
    % 修复：调整行高缩放的正确语法
    \renewcommand{\arraystretch}{1.4} 
    
    % 2. 夹心层：圆角盒子 (tcolorbox config)
    \begin{tcolorbox}[
        enhanced,
        width=\textwidth,                       
        % 修复核心：使用 >{\raggedright\arraybackslash}X 让文字靠左对齐，取消两端对齐的硬拉伸
        tabularx={
            >{\raggedright\arraybackslash}X | 
            >{\raggedright\arraybackslash}X | 
            >{\raggedright\arraybackslash}X | 
            >{\raggedright\arraybackslash}X | 
            >{\raggedright\arraybackslash}X | 
            >{\raggedright\arraybackslash}X
        },
        arc=8pt,                                
        boxrule=0.1pt,                            
        colback=white,
        colframe=tableheader,
        clip upper,                             
        boxsep=0pt, left=0pt, right=0pt, top=0pt, bottom=0pt 
    ]
        % 3. 内层：直接写表格的数据内容
        \rowcolor{tableheader} 
        % 表头建议保持居中对齐，所以这里局部加上 \centering
        \multicolumn{1}{c|}{\textbf{\color{white}实现方式}} & 
        \multicolumn{1}{c|}{\textbf{\color{white}性能体验}} & 
        \multicolumn{1}{c|}{\textbf{\color{white}显卡/GPU 支持}} & 
        \multicolumn{1}{c|}{\textbf{\color{white}系统隔离度}} & 
        \multicolumn{1}{c|}{\textbf{\color{white}切换便利性}} & 
        \multicolumn{1}{c}{\textbf{\color{white}适合人群}} \\ \hline
        
        \texttt{虚拟机} & 较差，资源开销大 & 极难配置，性能极差 & 完全隔离，最安全 & 窗口化运行，极方便 & 偶尔体验、测试轻量软件 \\ \hline
        
        \texttt{WSL2} & 接近原生，CPU/内存效率高 & 支持 CUDA，但有虚拟化损耗 & 与 Windows 高度融合 & 终端秒开，极方便 & Windows 用户，兼顾普通开发 \\ \hline
        
        \texttt{双系统} & 完美，100\% 原生硬件性能 & 原生支持，完美释放 GPU 算力 & 完全独立，互不干扰 & 必须重启电脑才能切换 & 硬核开发者、Linux 主力用户 \\
    \end{tcolorbox}
\end{table}

\section{环境搭建}
在运行比赛之前，我们需要安装 \texttt{rcssserver3d} 和相应的依赖。

\begin{tipbox}[Ubuntu 版本选择]
强烈推荐使用 Ubuntu 20.04 或 22.04。在虚拟机中运行 Server 可能会因为 OpenGL 支持不足导致监视器 (Monitor) 崩溃，建议双系统开发。
\end{tipbox}

\subsection{安装依赖}
以下是安装基础依赖的命令：
\begin{lstlisting}[language=bash]
sudo apt-get update
sudo apt-get install build-essential cmake libboost-all-dev
\end{lstlisting}

\section{智能体底层架构}
RoboCup 3D 的智能体主要通过 TCP 协议与 Server 通信。以 UTAustinVilla 代码库为例，其核心分为 \textbf{感知(Sense)} 和 \textbf{行动(Act)}。

\subsection{逆运动学 (Inverse Kinematics)}
为了让 Nao 机器人走路，我们需要计算关节角度。以下是计算目标末端执行器位置的基础公式：
\begin{equation}
    T_e = \prod_{i=1}^{n} A_i(\theta_i)
\end{equation}
其中 $\theta_i$ 是第 $i$ 个关节的角度。

\section{常见问题 (FAQ)}
\begin{itemize}
    \item \textbf{智能体连不上服务器？} 检查是否修改了 \texttt{rcssserver3d} 默认的 3100 端口。
    \item \textbf{机器人一直摔倒？} 检查 ZMP (零力矩点) 的计算是否在支撑多边形内。
\end{itemize}

\end{document}
